\section{Funtores sobre una categoría con clivaje}
\label{VI.12}

Sea $\cal{F}$ una categoría con clivaje normalizado sobre~$\cal{E}$.
Para todo objeto $S$ de~$\cal{E}$, denotemos por
$$
i_S\colon \cal{F}_S\to\cal{F}
$$
el funtor de inclusión. Se tiene, pues, una transformación natural, para
todo morfismo $f\colon T\to S$ de~$\cal{E}$:
$$
\alpha_f\colon i_T f^*\to i_S,
$$
donde $f^*$ es el funtor de cambio de base
$\cal{F}_S\to\cal{F}_T$ para $f$ definido por el clivaje. Sea ahora
$$
F\colon\cal{F}\to\cal{C}
$$
un funtor de~$\cal{F}$ en una categoría $\cal{C}$. Pongamos, para todo
$S\in\Ob(\cal{E})$,
$$
F_S=F\circ i_S\colon\cal{F}_S\to\cal{C}
$$
y, para todo $f\colon T\to S$ de~$\cal{E}$,
$$
\varphi_f=F*\alpha_f\colon F_T f^*\to F_S.
$$
Así, a todo funtor $F\colon\cal{F}\to\cal{C}$ se le asocian una familia
$(F_S)$ de funtores $\cal{F}_S\to\cal{C}$ y una familia $(\varphi_f)$ de
transformaciones naturales $F_T f^*\to F_S$. Estas familias satisfacen
las condiciones siguientes:

a) $\varphi_{\id_S}=\id_{F_S}$.

b) Para dos morfismos $f\colon T\to S$ y $g\colon U\to T$ de~$\cal{E}$,
es conmutativo el cuadrado de transformaciones naturales
$$
\xymatrix@C=2cm{ F_U g^* f^* \ar[r]^-{F_U*c_{f,g}} \ar[d]_{\varphi_g * f^*} &
F_U(fg)^* \ar[d]^{\varphi_{fg}} \\ F_T f^*\ar[r]^-{\varphi_f} & F_S.}
$$
\marginpar{190}%
La primera relación es inmediata, y la segunda se obtiene aplicando el
funtor~$F$ al diagrama conmutativo
$$
\xymatrix@C=2cm{ g^* f^* (\xi) \ar[r]^-{c_{f,g}(\xi)}
\ar[d]_{\alpha_g(f^*(\xi))} & (fg)^*(\xi) \ar[d]^{\alpha_{fg}(\xi)}\\
f^*(\xi) \ar[r]^-{\alpha_f(\xi)}& \xi }
$$
para un objeto variable $\xi$ de~$\cal{F}_S$.

Si $G$ es un segundo funtor $\cal{F}\to\cal{C}$, que da lugar a funtores
$G_S\colon\cal{F}_S\to\cal{C}$ y a transformaciones naturales
\ifthenelse{\boolean{orig}}
{$\psi_f\colon G_Tf\to G_S$}
{$\psi_f\colon G_Tf^*\to G_S$},
y si $u\colon F\to G$ es una transformación natural, le corresponden las
transformaciones naturales $u*i_S$:
$$
u_S\colon F_S\to G_S,
$$
y se comprueba de inmediato que, para todo morfismo $f\colon T\to S$ de
$\cal{E}$, son conmutativos los cuadrados
$$
\begin{array}{c}
\xymatrix@C=1.5cm{ F_Tf^* \ar[r]^-
{\ifthenelse{\boolean{orig}}{\phi_f}{\varphi_f}}
\ar[d]_{u_T * f^*} & F_S
\ar[d]^{u_S} \\ G_T f^* \ar[r]^-{\psi_f}& G_S}
\end{array}
\leqno\quad \textup{c)}
$$

\begin{proposition}
\label{VI.12.1}
Sea $\cal{H}(\cal{F},\cal{C})$ la categoría cuyos objetos son los pares
formados por familias $(F_S)$
\ifthenelse{\boolean{orig}}
{($S\in\Ob(\cal{F})$)}
{($S\in\Ob(\cal{E})$)}
de funtores $\cal{F}_S\to\cal{C}$ y familias $(\varphi_f)$
\ifthenelse{\boolean{orig}}
{($f\in\Fl(\cal{F})$)}
{($f\in\Fl(\cal{E})$)}
de transformaciones naturales $F_Tf^*\to F_S$ que satisfacen las
condiciones \emph{a) y~b)}, y cuyos morfismos son las familias $(u_S)$
\ifthenelse{\boolean{orig}}
{($S\in\Ob(\cal{F})$)}
{($S\in\Ob(\cal{E})$)}
de transformaciones naturales
\marginpar{191}
$F_S\to G_S$ que verifican la condición de conmutatividad \emph{c)}
escrita más arriba (la composición de los morfismos se efectúa componiendo
las transformaciones naturales entre funtores $\cal{F}_S\to\cal{C}$).
Entonces las dos reglas explicitadas más arriba definen un
\emph{isomorfismo} $K$ de la categoría
$\SheafHom(\cal{F},\cal{C})$ con la categoría
$\cal{H}(\cal{F},\cal{C})$.
\end{proposition}

Es inmediato que se tiene así un \emph{funtor} de la primera categoría en
la segunda. Este funtor es plenamente fiel, pues, dados $F,G$, la
aplicación $\Hom(F,G)\to\allowbreak\Hom(K(F),K(G))$ es manifiestamente
inyectiva. Para demostrar que es sobreyectiva, basta observar que la
condición de conmutatividad c) expresa la functorialidad de las
aplicaciones
$u(\xi)=u_S(\xi)\colon
\ifthenelse{\boolean{orig}}
{F(S)}
{F(\xi)}
=F_S(\xi)\to G(\xi)=G_S(\xi)$ para los morfismos de la forma
$\alpha_f(\xi)$ en~$\cal{F}$. Por otra parte, se tiene la functorialidad
en cada categoría fibra, es decir, para los morfismos de~$\cal{F}$ que
son $T$-morfismos ($T\in\Ob(\cal{E})$), de donde resulta la functorialidad
para todo morfismo de~$\cal{F}$,
\ifthenelse{\boolean{orig}}
{puesto que un $f$-morfismo (donde $f\colon T\to S$ es un morfismo en
($\cal{E}$) es de manera única) un compuesto de un morfismo
$\alpha_f(\xi)$ y de un $T$-morfismo.}
{puesto que todo $f$-morfismo (donde $f\colon T\to S$ es un morfismo en
$\cal{E}$) es, de manera única, un compuesto de un morfismo
$\alpha_f(\xi)$ y de un $T$-morfismo.}
Queda por demostrar que el funtor~$K$ es biyectivo sobre los objetos. El
argumento precedente muestra ya que $K$ es inyectivo sobre los objetos;
queda probar que es sobreyectivo, es decir, que, si se parte de un sistema
$(F_S),(\varphi_f)$ que satisface a) y b), y se define una aplicación
$\Ob\cal{F}\to\Ob\cal{C}$ por
$$
F(\xi)=F_S(\xi) \quad\text{para }
\xi\in\Ob\cal{F}_S\subset\Ob\cal{F}
$$
y una aplicación $\Fl(\cal{F})\to\Fl(\cal{C})$ por
$$
F(\alpha_f(\xi)u')=\varphi_f(\xi)\,F_T(u'),
$$
para todo morfismo $f\colon T\to S$ de~$\cal{E}$, todo objeto $\xi$
de~$\cal{F}_S$ y todo $T$-morfismo $u'$ con destino $f^*(\xi)$, se
obtiene un \emph{funtor} $F$ de~$\cal{F}$ en~$\cal{C}$. En efecto, la
relación $F(\id_\xi)=\id_{F(\xi)}$ es inmediata; queda demostrar la
multiplicatividad $F(uv)=F(u)F(v)$ cuando se tiene un $f$-morfismo
$u\colon\eta\to\xi$ y un $g$-morfismo
\ifthenelse{\boolean{orig}}
{$v\colon\zeta\to\nu$}
{$v\colon\zeta\to\eta$},
con $f\colon T\to S$ y $g\colon U\to T$ morfismos de~$\cal{E}$.
Poniendo $w=uv$, se tendrá
$$
u=\alpha_f(\xi)u'\quoi,\quad v=\alpha_g(\eta)v'\quoi,\quad
w=\alpha_{fg}(\xi)w'.
$$
\marginpar{192}%
con
$$
w'=c_{f,g}(\xi)g^*(u')v'\quad \hbox{(\cf \No \Ref{VI.8})}.
$$

Con esta notación, hay que demostrar la conmutatividad del contorno
exterior del diagrama siguiente:
$$
\xymatrix{ F_U(\zeta) \ar@/^1.5pc/[rrrrrr]^{F(w')}
\ar[rr]^{\ifthenelse{\boolean{orig}}{F_u(v')}{F_U(v')}}
\ar[drr]_{F(v)} & & F_Ug^*(\eta)
\ar[rr]^{\ifthenelse{\boolean{orig}}{F_ug^*(u')}{F_Ug^*(u')}}
\ar[d]^{\varphi_g(\eta)} & & F_Ug^*f^*(\xi)
\ar[rr]^{F_U(c_{f,g}(\xi))}\ar[d]^{\varphi_g(f^*(\xi))}& &
F_U(fg)^*(\xi) \ar[d]^{\varphi_{fg}(\xi)} \\ & & F_T(\eta)
\ar@/_1.2pc/[rrrr]_{F(u)} \ar[rr]_{F_T(u')} & & F_Tf^*(\xi)
\ar[rr]_{\varphi_f(\xi)} & & F_S(\xi) }
$$
Ahora bien, el triángulo izquierdo es conmutativo por definición de
$F(v)$; el cuadrado central es conmutativo porque se deduce del morfismo
\ifthenelse{\boolean{orig}}
{$u'\colon \xi\to f(\eta)$}
{$u'\colon \eta\to f^*(\xi)$}
mediante la transformación natural $\alpha_g$;
por último, el cuadrado derecho es conmutativo en virtud de la condición
b). De ello resulta la conclusión buscada.

Supongamos ahora que $\cal{C}$ es también una categoría con clivaje
normalizado sobre~$\cal{E}$, que en adelante llamaremos $\cal{G}$, y que
nos interesan los $\cal{E}$-funtores de~$\cal{F}$ en~$\cal{G}$. Si $F$ es
uno de tales funtores, induce funtores
$$
F_S\colon \cal{F}_S\to\cal{G}_S
$$
entre las categorías fibra. Por otra parte, para todo morfismo
$f\colon T\to S$ de~$\cal{E}$ y todo objeto $\xi$ de~$\cal{F}_S$, el
$f$-morfismo $F(\alpha_f(\xi))$ se factoriza de manera única por un
$T$-morfismo
$$
\varphi_f(\xi)\colon F_T(f^*_{\cal{F}}(\xi))\to
f^*_{\cal{G}}(F_S(\xi))
$$
(donde el subíndice $\cal{F}$ o $\cal{G}$ indica la categoría con
clivaje para la que se toma el funtor de imagen inversa), de donde se
obtiene una transformación natural entre funtores de~$\cal{F}_S$ en
$\cal{G}_T$:
$$
\varphi_f\colon F_T f^*_{\cal{F}}\to f^*_{\cal{G}}F_S.
$$
\marginpar{193}%
Los dos sistemas $(F_S)$ y $(\varphi_f)$ satisfacen las condiciones
siguientes:

a') $\varphi_{\id_S}=\id_{F_S}$.

b') Para dos morfismos $f\colon T\to S$ y $g\colon U\to T$ de
$\cal{E}$, es conmutativo el diagrama siguiente de transformaciones
naturales:
$$
\xymatrix{ F_U
\ifthenelse{\boolean{orig}}{g_F^*f_F^*}{g_{\cal{F}}^*f_{\cal{F}}^*}
\ar[rr]^{F_U *c_{f,g}^\cal{F}}
\ar[d]_{\ifthenelse{\boolean{orig}}{\varphi_{g^*}f^*_{\cal{F}}}{\varphi_g*f^*_{\cal{F}}}} & &
F_U\ifthenelse{\boolean{orig}}{(fg)^*_F}{(fg)^*_{\cal{F}}}
\ar[dd]^{\varphi_{fg}} \\ g^*_{\cal{G}}F_T
\ifthenelse{\boolean{orig}}{f^*_F}{f^*_{\cal{F}}}
\ar[d]_{g^*_{\cal{G}}*\varphi_f} & & \\ g^*_{\cal{G}}f^*_{\cal{G}}F_S
\ar[rr]^{c_{f,g}^\cal{G}*F_S} & & (fg)^*_{\cal{G}}F_S. }
$$
Dejamos al lector su verificación, así como el enunciado y la demostración
del análogo de la Proposición~\Ref{VI.12.1}; este implica que se obtiene
así una correspondencia biunívoca entre el conjunto de los
$\cal{E}$-funtores de~$\cal{F}$ en~$\cal{G}$ y el conjunto de los
sistemas $(F_S)$,$(\varphi_f)$ que satisfacen las condiciones a') y b')
anteriores. Por supuesto, en esta correspondencia, los funtores
cartesianos se caracterizan por la propiedad de que las transformaciones
$\varphi_f$ son isomorfismos.

\begin{remarquestar}
Por supuesto, casi siempre conviene razonar directamente con categorías
fibradas sin utilizar clivajes explícitos, lo cual evita, en particular,
tener que recurrir, para la noción sencilla de $\cal{E}$-funtor o de
$\cal{E}$-funtor cartesiano, a una interpretación pesada como la anterior.
Para evitar pesadeces insoportables y obtener enunciados más intrínsecos,
\marginpar{194}
hemos tenido que renunciar a partir, como en~\cite{VI.2}, de la noción de
categoría con clivaje (llamada «categoría fibrada» en \loccit), que pasa a
segundo plano en favor de la de categoría fibrada. Es, por lo demás,
probable que, contrariamente al uso todavía predominante, ligado a antiguos
hábitos de pensamiento, acabe por resultar más cómodo, en los problemas
universales, no poner el acento en \emph{una} solución supuestamente elegida
de una vez por todas, sino situar todas las soluciones en pie de igualdad.
\end{remarquestar}

\begin{sgabibliography}{0}{VI.13}
\bibitem[1]{VI.1} A\ptbl \textsc{Grothendieck}, Sur quelques points d'algèbre
homologique,
\ifthenelse{\boolean{orig}}
{Tohoku Math.\ Journal, pp\ptbl 119--221, Vol\ptbl 9, 1957.}
{T\^ohoku Math.\ J. \textbf{9} (1957) p\ptbl 119--221.}

\bibitem[2]{VI.2} A\ptbl \textsc{Grothendieck}, Technique de descente et
Théorèmes d'existence,~I, Séminaire Bourbaki~190, diciembre de~1959.
\end{sgabibliography}


\refstepcounter{chapter}
\addtocontents{toc}{\protect\medskip\protect\par\protect\noindent\protect\textbf{VII: no existe}}
